\begin{theorem}[Proved: the positive cardinal-filtration equality]\label{cat:I11}
Let $E=\operatorname{Ind}_\kappa(C)$, let $\mu\geq\kappa$ be infinite,
and let $\operatorname{Fil}^{\mu}_C(E)$ consist of the objects admitting
a $\kappa$-directed presentation by objects of $C$ with a cofinal index
set of cardinal at most $\mu$.  Always
\[
 \operatorname{Fil}^{\mu}_C(E)\subseteq E_{\mu^+}.
\]
If $\kappa\triangleleft\mu^+$, equivalently
$\operatorname{cf}([\mu]^{<\kappa},\subseteq)\leq\mu$, then equality
holds.  In particular $\mu^{<\kappa}=\mu$ suffices.  This proves exactly
the positive assertion on 161-2/060(b): with its convention
$\kappa=\pi^+$, every $\mu=2^c$ with $c\geq\pi$ satisfies the hypothesis.
At these ranks $E$ is $\mu^+$-accessible and
$E\simeq\operatorname{Ind}_{\mu^+}(E_{\mu^+})$.
If $E$ is cocomplete, the latter accessibility and Ind equivalence hold
at every regular $\lambda\geq\kappa$; this does not force the equality
of the two filtrations at every rank, as Proposition~\ref{cat:I10} shows.
\end{theorem}
\begin{proof}
A cofinal subset of a directed poset is still $\kappa$-directed.  Thus
an object of $\operatorname{Fil}^{\mu}_C(E)$ has a presentation with at
most $\mu$ objects and arrows.  Each object of $C$ is
$\mu^+$-presentable, and a colimit of at most $\mu$ such objects and
arrows remains $\mu^+$-presentable: its Hom functor is the corresponding
limit of at most $\mu$ set-valued functors.  Such limits commute with
$\mu^+$-filtered colimits of sets.  This proves the inclusion without
a cardinal-arithmetic hypothesis.

Under the sharp-cardinal hypothesis, Theorem~\ref{cat:I9}'s closure
property permits the construction of $J_\mu$ and the decomposition in
Lemmas~\ref{cat:I2}--\ref{cat:I3} with precisely the same bound $\mu$.
Consequently every $X\in E_{\mu^+}$ is a retract of some
$Y=\operatorname{colim}_{j\in J}c_j$, where $J$ is $\kappa$-directed
and $|J|\leq\mu$.  Indeed $X$ is a $\mu^+$-directed colimit of these
objects, so its identity factors through one of them.
It remains to remove this retract; it must not simply be omitted.

Write $i:X\to Y$, $p:Y\to X$, $pi=1$, $e=ip$, and let
$a_j:c_j\to Y$ be the presentation maps.  By
$\kappa$-presentability, for each $j$ there are $k(j)\in J$ and
$f_j:c_j\to c_{k(j)}$ satisfying $a_{k(j)}f_j=ea_j$.
Put $q_j=pa_j$.  Inside $C/X$ take the subcategory $D$ generated by
the presentation arrows and these additional arrows $f_j$; the latter
belong to $C/X$ because $q_{k(j)}f_j=q_j$.  There are at most $\mu$
objects and arrows after closing under finite composition.  Its
tautological diagram has colimit $X$: a cocone on the original
presentation is a map $g:Y\to Z$, and the additional conditions are
exactly $ge=g$.  Such maps are uniquely $hp$ for $h:X\to Z$.

We enlarge $D$ to a $\kappa$-filtered subcategory $D'\subset C/X$
with at most $\mu$ arrows, preserving this colimit.  At each of
$\kappa$ stages choose a cofinal family, of cardinal at most $\mu$,
of the subsets of the current arrow set having size less than
$\kappa$.  For each chosen subset add a cocone in $C/X$ on its
underlying graph.  Such a cocone exists: first dominate its fewer than
$\kappa$ objects, then equalize its fewer than $\kappa$ arrow
conditions, using $\kappa$-filteredness of $C/X$.
For every new object $(c,q)$ also add one arrow from it to an original
object $(c_j,q_j)$.  This is possible because $iq:c\to Y$ factors
through some $a_j$, and applying $p$ gives the required factorization
of $q$.  Close under composition and take unions at limit stages.
All stages and their union have at most $\mu$ arrows.
Regularity of $\kappa$ ensures that every diagram with fewer than
$\kappa$ arrows in the union occurs at one stage and acquires a cocone
at the next.  Hence $D'$ is $\kappa$-filtered.
Its colimit remains $X$: restriction of any cocone to $D$ gives its
unique map $X\to Z$, and the chosen outgoing arrows to $D$ force
the component at every added object to be the corresponding composite.

Finally the directed-poset replacement of Lemma~\ref{cat:I7} can be
made with at most $\mu$ vertices.  In that construction choose a
cofinal family of fewer-than-$\kappa$ subsets at each stage instead
of all such subsets; there are at most $\mu$ choices by the same
sharp-cardinal hypothesis.  The additional pairs used for comma-category
connectivity are also at most $\mu$, and there are $\kappa\leq\mu$
stages.  This yields a $\kappa$-directed presentation of $X$ of size
at most $\mu$, proving equality.

There is only a set of diagrams of size at most $\mu$ in the small
category $C$.  The decomposition by $J_\mu$ therefore proves
$\mu^+$-accessibility, and Theorem~\ref{cat:I1} gives the displayed
Ind equivalence.  If $\mu=2^c$ and $c\geq\pi$, then for every
$\alpha<\kappa=\pi^+$ one has
$(2^c)^\alpha=2^{c\alpha}=2^c$ when $\alpha>0$; hence
$\mu^{<\kappa}=\mu$.
For the final assertion, a small dense category of
$\kappa$-presentables is still dense and consists of
$\lambda$-presentables.  The cocomplete clause of
Theorem~\ref{cat:I1} applies at every such regular $\lambda$.
\end{proof}

